```latex
\section*{Computation Evidence}

\textbf{Statement.} It was verified that for every natural number $n \leq 2{,}540{,}160$, the equation $n = \sum_{d \in \text{digits}(n)} d!$ holds if and only if $n \in \{1, 2, 145, 40585\}$.

\textbf{Method.} The verification was carried out by exhaustive enumeration of all natural numbers up to $2{,}540{,}160$, distributed across 3 cluster nodes (node1, node2, node3), each scanning a chunk of $846{,}720$ integers and computing the digit-factorial sum for every candidate. The complete range was partitioned into 3 non-overlapping chunks covering $[0, 2{,}540{,}160]$ with no gaps or overlaps.

\textbf{Evidence Summary.} Across all 3 chunks totalling $2{,}540{,}160$ integers scanned, exactly 4 fixed points of the digit-factorial sum were found: $n \in \{1, 2, 145, 40585\}$; all remaining values satisfied $n \neq \sum_{d} d!$. The empirical verdict was deemed \emph{sufficient}, confirming that no factorions exist in the range $[0, 2{,}540{,}160]$ beyond the four known values.
```